% Temperature: the nonlinearity N = ALL - GHG - AAER as a map, 2031-2050.
% Built from the gridded NetCDFs by scripts/make_fig1112_maps.py
% — do not edit by hand. The caption below says what the numbers
% are; \input this file directly.
\begin{table}[htbp]
\centering
\footnotesize
\setlength{\tabcolsep}{4pt}
\caption{The nonlinearity $N = \Delta(\mathrm{ALL}) - \Delta(\mathrm{GHG}) - \Delta(\mathrm{AAER})$ in temperature, averaged over 2031--2050 --- the last 20 years the single-forcing runs cover --- as area-weighted means over named regions. $\Delta$ is the anomaly relative to each side's own 1850--1900 mean, and $N$ is computed separately on each side, so neither is a reference for the other. Units are $^{\circ}$C; precipitation is left in mm\,day$^{-1}$ rather than converted to a percentage because a percentage residual divides by a near-zero climatology over the subtropical deserts. The area-weighted pattern correlation between the two maps is $r = 0.935$: the emulator reproduces where the nonlinearity is, and the ``difference'' column says by how much it overstates it. ``Significant'' is the percentage of the region's area where the two differ at $p < 0.05$, point-wise, with the false-discovery rate controlled at $q = 0.1$ across the whole map --- so MORE of it means a stronger detection of disagreement, not a better emulator. NOTE THAT $N$ IS NOT A PURE INTERACTION TERM: the single-forcing runs hold ozone, land use, biomass burning, solar and volcanic forcing at 1850 levels, so $N$ = (those forcings) + (nonlinearity), and the Amazon and Congo signals in particular are land use and biomass burning.}
\label{tab:fig11}
\begin{tabular}{|l|r|r|r|r|}
\hline
\textbf{Region} & \textbf{CESM2} & \textbf{Emulator} & \textbf{Difference} & \textbf{Significant} \\
 & \multicolumn{3}{c|}{($^{\circ}$C)} & (\%) \\
\hline
Global & +0.469 & +0.597 & +0.128 & 18 \\
\hline
Arctic & +1.482 & +1.867 & +0.385 & 1 \\
N. Atlantic & +1.564 & +2.058 & +0.494 & 21 \\
N. America & +0.998 & +1.089 & +0.091 & 5 \\
Europe & +1.059 & +1.196 & +0.137 & 13 \\
East Asia & +0.786 & +0.990 & +0.203 & 24 \\
South Asia & +0.371 & +0.424 & +0.052 & 24 \\
Sahel & +0.622 & +0.986 & +0.363 & 80 \\
Tropics & +0.396 & +0.424 & +0.028 & 17 \\
Antarctic & +0.099 & +0.106 & +0.006 & 6 \\
\hline
\end{tabular}
\end{table}
